% Détection de petits objets vus de drone
\section{Détection de petits objets vus de drone}
\label{sec:detection}

\subsection{La recette dominante : tête P2, attention, suppression de P5}

Les YOLO spécialisés drone du corpus modifient tous la pyramide de caractéristiques pour
les petits objets. Trois gestes reviennent : ajouter une tête P2 (stride 4, carte de
160$\times$160 pour une entrée de 640), ajouter de l'attention ou des blocs Transformer,
et supprimer la tête P5, qui porte beaucoup de paramètres mais sert surtout aux grands
objets.

\paragraph{TPH-YOLOv5 et TPH-YOLOv5++.} TPH-YOLOv5
\citep{arxiv-2108.11539-tph-yolov5-improved-yolov5-transformer}\pp{} est la référence de
cette famille : YOLOv5x, quatrième tête P2, têtes à Transformer, CBAM, test multi-échelle
et ensemble de cinq modèles. L'ensemble atteint 39,18 de mAP [.5:.95] sur le
test-challenge VisDrone 2021 \src{§4.2}\med{}, sur RTX 3090, sans aucune vitesse rapportée.
La tête P2 fait passer YOLOv5x de 219 à 259\,GFLOPs \src{§4.3}. \citet{2023-zhao-tph-yolov5-boosting-object-detection}
mesurent ce coût : la tête P2 seule apporte +2,1 points d'AP, fait passer la mémoire GPU de
4\,279 à 4\,667\,Mo et le débit de 13,68 à 10,89\,FPS \src{Table~5}. TPH-YOLOv5++ remplace
la tête P2 par un transfert d'information vers P3 (CA-Trans) et remonte à 11,86\,FPS pour
33,1 d'AP sur test-dev, contre 7,36\,FPS et 34,0 pour TPH-YOLOv5 \src{Table~5} ; toujours
sur RTX 3090 Ti, à 1536$\times$1536.

\paragraph{Variantes légères.} Le tableau~\ref{tab:yolo-drone} rassemble les autres
variantes. Toutes allègent le réseau en paramètres ; aucune ne mesure son coût sur une
cible embarquée, sauf Drone-YOLO.

\begin{table}[!htbp]
\centering\footnotesize
\caption{YOLO spécialisés drone du corpus, sur VisDrone. Plateformes de bureau sauf
mention. Protocoles hétérogènes (split, résolution, métrique) : la colonne précision ne
permet pas de classer.}
\label{tab:yolo-drone}
\begin{tabularx}{\textwidth}{@{}P{3.3cm}LP{2.0cm}P{2.6cm}@{}}
\toprule
Modèle & Modification principale & Précision & Coût, plateforme \\
\midrule
GCL-YOLO-S \citep{2023-cao-gcl-yolo-ghostconv-lightweight-yolo} & GhostConv, P2 ajoutée, P5
  retirée ; 1,64\,M param. & 39,6 mAP50 \src{Table~3} & 53\,FPS, GPU n/r \\
Efficient YOLOv7-Drone \citep{2023-fu-efficient-yolov7-drone-enhanced-object} & P2, P5 retirée,
  convolution creuse guidée par le premier plan & 32,1 AP (val, 640) \src{Table~2} & 89\,FPS
  contre 213 pour YOLOv7, RTX 3090 \src{Table~5} \\
Drone-YOLO nano \citep{2023-zhang-drone-yolo-efficient-neural-network} & P2, sandwich-fusion,
  RepVGG ; 3,05\,M param. & 38,1 mAP50 (val) \src{Table~5} & 127,4\,ms sur Jetson TX2
  \src{Table~6} \\
SF-YOLOv5 \citep{2022-liu-sf-yolov5-lightweight-small-object} & P5 retirée, PB-FPN ;
  2,23\,M param. & 34,3 mAP50\med{} \src{Table~4} & 11,5\,ms, RTX 2060 \\
DroneScan-YOLO \citep{arxiv-2604.13278-dronescan-yolo-redundancy-aware-lightweight}\pp &
  entrée 1280, branche P2, pruning non structuré & 55,3 mAP50 (val) \src{Table~4} & 96,7\,FPS,
  RTX 4090 portable \\
YOLO-Drone \citep{arxiv-2511.10905-yolo-drone-efficient-object-detection} & GhostConv dans la
  tête de YOLOv11n & 30,4 mAP50 \src{Table~2} & 1,8\,ms, Tesla T4 \\
DEAL-YOLO \citep{arxiv-2503.04698-deal-yolo-drone-efficient-animal}\pp & P2 à la place de P5,
  NWD ; faune (hors VisDrone) & 48,5 mAP50, BuckTales \src{Table~1} & \nr \\
\bottomrule
\end{tabularx}
\end{table}

\paragraph{Ce que coûte la tête P2.} Trois articles donnent un chiffre. Drone-YOLO nano
met 127,4\,ms par image sur Jetson TX2, contre 30,8\,ms pour YOLOv8n à nombre de
paramètres égal, soit quatre fois plus lent pour +4,1 points de mAP50
\citep[Table~6]{2023-zhang-drone-yolo-efficient-neural-network} ; l'auteur conclut
pourtant à des résultats « highly accurate » sans commenter la latence. Pour
Efficient YOLOv7-Drone, ajouter P2 fait tomber le débit de 212,8 à 137,0\,FPS sur RTX 3090
\citep[Table~5]{2023-fu-efficient-yolov7-drone-enhanced-object}. GCL-YOLO-S perd 7\,FPS
face à YOLOv5-S alors que ses GFLOPs baissent \citep[Table~3,
§5]{2023-cao-gcl-yolo-ghostconv-lightweight-yolo}. Le coût de P2 est donc un coût
d'activations et d'accès mémoire, que les FLOPs et le nombre de paramètres ne voient pas
\deduit. Sur un FPGA, où les activations d'une carte 160$\times$160 doivent tenir en
mémoire on-chip, c'est précisément ce coût qui domine \deduit.

\paragraph{Paramètres et FLOPs ne prédisent pas la latence.} Le constat revient dans cinq
articles. SF-YOLOv5 retire 68\,\% des paramètres pour 7\,\% de latence en moins
\citep[Table~3]{2022-liu-sf-yolov5-lightweight-small-object}. SlimYOLOv3
\citep{arxiv-1907.11093-slimyolov3-narrower-faster-better-real}\pp{} (version ICCVW~2019
probable) retire 90,8\,\% des FLOPs de YOLOv3-SPP3 par élagage de canaux pour 49,5\,\% de
temps en moins sur GTX 1080 Ti \src{§5}, et reste plus lent que YOLOv3-tiny à FLOPs
comparables. F-YOLOv7 divise les FLOPs par 2,8 sans gain de débit sur RTX 4090
\citep[Table~2]{2025-du-f-yolov7-fast-robust-real}. DroneScan-YOLO reconnaît que son
pruning non structuré ne retire aucun paramètre et n'accélère rien sur matériel embarqué
\citep[§5]{arxiv-2604.13278-dronescan-yolo-redundancy-aware-lightweight}. YOLO-GCOF,
annoncé pour Raspberry Pi 4B, ne donne que des GFLOPs
\citep[§IV-D8]{2025-yu-yolo-gcof-lightweight-low-altitude}. Tous ces résultats sont sur
GPU de bureau : sur FPGA, NPU ou MCU, l'effet peut s'inverser (GhostConv et depthwise
sont souvent mal servis par un DPU) \deduit.

\subsection{Calculer seulement où sont les objets}

Une deuxième famille garde un détecteur à résolution modeste et concentre le calcul sur
les régions utiles.
\begin{itemize}
  \item \textbf{Tuilage} : SAHI \citep{arxiv-2202.06934-slicing-aided-hyper-inference-fine}\pp{}
    (publié à ICIP~2022 d'après son DOI) découpe l'image en tuiles 640$\times$640. Avec
    fine-tuning sur tuiles, TOOD passe de 29,4 à 43,5 d'AP50 sur VisDrone test-dev, et
    son AP50$_S$ de 18,1 à 31,7 \src{Table~1}. Le coût, linéaire en nombre de tuiles
    \src{§5}, n'est pas mesuré.
  \item \textbf{Régions denses} : ClusDet
    \citep{arxiv-1904.08008-clustered-object-detection-aerial-images}\pp{} détecte des
    groupes d'objets puis les traite à part ; il atteint 26,7 d'AP et 17,6 d'AP$_S$ en R-50
    \src{Table~1}, mais met 0,273\,s par image sur GTX 1080 Ti, cinq fois plus que la
    baseline sans découpage. UFPMP-Det
    \citep{arxiv-2112.10415-ufpmp-det-accurate-efficient-object}\pp{} assemble les régions
    en une mosaïque unique et atteint 36,6 d'AP en R-50 sur VisDrone val, pour 0,152\,s par
    image \src{Tables~1, 3}.
  \item \textbf{Requêtes éparses} : QueryDet
    \citep{arxiv-2103.09136-querydet-cascaded-sparse-query-accelerating}\pp{} ne calcule
    la tête haute résolution qu'aux positions désignées par une carte basse résolution. Sur
    COCO, il garde le gain d'AP$_S$ de P2 (24,33 contre 22,64) à 14,88\,FPS, contre 4,85
    pour une tête P2 dense \src{Table~1}. Sur VisDrone, il plafonne à 2,75\,FPS
    \src{Table~2}.
  \item \textbf{Sortie anticipée} : BlankSkip
    \citep{arxiv-2603.28149-blankskip-early-exit-object-detection}\pp{} saute le reste du
    réseau sur les images vides ; il est traité avec les nano-drones
    (section~\ref{sec:nano}).
\end{itemize}
Aucune de ces méthodes n'est mesurée sur cible embarquée, sauf BlankSkip. Aucune ne
rapporte la latence au pire cas ni la gigue, alors que leur coût dépend du contenu de la
scène. Du point de vue matériel, elles ne se valent pas : la mosaïque de taille fixe
d'UFPMP-Det ou le tuilage de SAHI gardent un calcul dense et régulier, favorable à un
accélérateur à entrée fixe, alors que la convolution éparse de QueryDet repose sur des
accès irréguliers \deduit.

\subsection{Protocoles et comparabilité}

Les chiffres VisDrone du corpus ne se comparent presque jamais entre eux :
\begin{itemize}
  \item le split varie (validation pour Efficient YOLOv7-Drone, ClusDet, UFPMP-Det,
    DroneScan-YOLO ; test-dev pour TPH-YOLOv5 et SAHI ; non précisé pour GCL-YOLO et
    YOLO-Drone) ;
  \item la métrique varie (mAP50 pour les YOLO, AP COCO pour les autres) ;
  \item la résolution varie de 640 à 1920\,px, et DroneScan-YOLO compare son modèle à
    1280\,px à des baselines à 640\,px \citep[§5]{arxiv-2604.13278-dronescan-yolo-redundancy-aware-lightweight} ;
  \item les baselines sont souvent réentraînées sans pré-entraînement, ou reprises des
    papiers d'origine.
\end{itemize}
Surtout, les gains annoncés sont souvent de l'ordre du point, sur un seul entraînement et
sans écart-type : +0,9 point pour \citet{2022-jung-improved-yolov5-efficient-object-detection},
+0,5 pour YOLO-Drone, +1,3 pour SF-YOLOv5. Enfin, l'AP$_S$ n'est rapportée que par trois
articles de ce groupe (ClusDet, SAHI, Efficient YOLOv7-Drone, ce dernier seulement pour ses
modèles intermédiaires : 19,4 à 22,0\,\% \citep[Tables~9, 12]{2023-fu-efficient-yolov7-drone-enhanced-object}),
alors que tous visent les petits objets.
